\documentclass[11pt]{article}

% ─── Packages ────────────────────────────────────────────────────────────────
\usepackage{amsmath,amssymb,amsthm}
\usepackage[margin=1in]{geometry}
\usepackage{hyperref}
\usepackage{graphicx}
\usepackage{booktabs}
\usepackage{array}
\usepackage{multirow}
\usepackage{xcolor}
\usepackage{authblk}
\usepackage{setspace}
\usepackage{caption}
\usepackage{subcaption}
\usepackage{enumitem}
\usepackage{microtype}
\usepackage{siunitx}

\hypersetup{
  colorlinks=true,
  linkcolor=blue,
  citecolor=blue,
  urlcolor=blue
}

% ─── Theorem environments ────────────────────────────────────────────────────
\theoremstyle{plain}
\newtheorem{theorem}{Theorem}[section]
\newtheorem{lemma}[theorem]{Lemma}
\newtheorem{proposition}[theorem]{Proposition}
\newtheorem{corollary}[theorem]{Corollary}
\theoremstyle{definition}
\newtheorem{definition}{Definition}[section]
\theoremstyle{remark}
\newtheorem{remark}{Remark}[section]

% ─── Title block ─────────────────────────────────────────────────────────────
\title{\textbf{Equal-Strength Donor Competition as a Mechanism\\
for Coordination Entropy Tuning:\\
LiFSI in DME/Trimethyl Phosphate as a Designed\\
Approach to SCE* = 1.466}}

\author[1]{Eric Robert Lawson}

\affil[1]{Independent Researcher, OrganismCore Initiative\\
\texttt{ORCID: 0009-0002-0414-6544}\\
Repository: \url{https://github.com/Eric-Robert-Lawson/attractor-oncology}}

\date{April 4, 2026\\
\small Pre-registration record:
\url{https://github.com/Eric-Robert-Lawson/attractor-oncology}\\
\small (timestamped commit history constitutes the priority
and pre-registration record)\\
\small Part of the SCE Framework preprint series.
Preprint~1 (foundation): same repository.}

% ─── Document ────────────────────────────────────────────────────────────────
\begin{document}

\maketitle

% ─── Abstract ────────────────────────────────────────────────────────────────
\begin{abstract}
Preprint~1 of this series derived the fixed point
$\mathrm{SCE}^{*} = 1.466$ as the Shannon Coordination Entropy of the
Li$^+$ first solvation shell that simultaneously maximises room-temperature
(RT) and low-temperature (LT) Coulombic efficiency (CE) in lithium metal
batteries.  The present paper introduces and formalises the
\textbf{ESP-matching mechanism} as the first rationally designed approach
to reaching that fixed point from below: when two co-solvents present
electrostatically equivalent binding sites to Li$^+$, the ion cannot
preferentially occupy either coordination class, the first-shell population
splits across both structural families, and SCE rises toward
$\mathrm{SCE}^{*}$.

We identify trimethyl phosphate (TMP) and 1,2-dimethoxyethane (DME) as the
canonical ESP-matched pair.  Computational electrostatic potential surface
analysis yields TMP ESP$_{\min} = -1.7402$~eV and DME ESP$_{\min} =
-1.7328$~eV, a difference of $0.007$~eV---within thermal noise at
operating temperature.  No prior work selects co-solvents by ESP matching
to the base solvent for the explicit purpose of Li$^+$ coordination entropy
tuning.  This is a new design principle.

Candidate~4 (LiFSI 1.0--1.2~M in DME/TMP, ratio scan targeting
$\mathrm{SCE} = 1.466$) and Candidate~4b (LiFSI in DME/fluorophosphonate,
weaker donor variant) are derived, with exact falsification conditions
pre-registered in the public repository before experiment.  A molecular
search across the phosphate ester class confirms TMP as the universal
cross-class attractor: every ether-anchored search in accessible chemical
space terminates at TMP or a close phosphate ester analogue.  The ESP
matching principle is shown to extend beyond the DME/TMP pair to any
binary system in which the solvent ESP minima are matched within
$\sim\!0.01$~eV, providing a general co-solvent selection rule for
coordination entropy engineering.
\end{abstract}

\vspace{1em}
\noindent\textbf{Keywords:} lithium metal battery, electrolyte design,
electrostatic potential, coordination entropy, solvation structure,
trimethyl phosphate, DME, co-solvent selection, fixed point, SCE*

\tableofcontents
\newpage

% ─── 1. Introduction ─────────────────────────────────────────────────────────
\section{Introduction}

\subsection{Context: The Fixed Point and the Design Problem}

Preprint~1 of this series established the following result: there exists a
unique value of the Shannon Coordination Entropy of the Li$^+$ first
solvation shell,

\begin{equation}
  \mathrm{SCE}^{*} = 1.466,
  \label{eq:fixed_point}
\end{equation}

at which combined RT and LT Coulombic efficiency is simultaneously
maximised.  The result is derived analytically from the intersection of
two opposing empirical gradients and confirmed across a 21-system published
dataset ($R^2 = 0.828$, $p = 4.5 \times 10^{-6}$).

The fixed point converts an unbounded experimental search into a bounded
navigation problem:

\begin{enumerate}[leftmargin=2em]
  \item Measure the Li$^+$ first-shell coordination geometry population.
  \item Compute $\mathrm{SCE}$ from the Shannon entropy of that
        distribution.
  \item If $\mathrm{SCE} < \mathrm{SCE}^{*}$: the shell is too ordered.
        Add structural diversity.
  \item If $\mathrm{SCE} > \mathrm{SCE}^{*}$: the shell is too disordered.
        Remove structural diversity.
  \item Iterate to target.
\end{enumerate}

The question that Preprint~1 leaves open---and that the present paper
addresses---is:

\begin{quote}
\itshape By what mechanism does a formulation chemist add structural
diversity to an ordered Li$^+$ solvation shell in a principled,
predictable, and rationally designed way?
\end{quote}

The answer introduced here is the \textbf{ESP-matching mechanism}.

\subsection{What Has Been Done Before}

Co-solvent selection in Li metal battery electrolyte design has historically
been guided by:

\begin{itemize}[leftmargin=2em]
  \item \textbf{Donor number (DN):} the enthalpy of coordination of the
        solvent to a reference Lewis acid.  Solvents with high DN
        coordinate Li$^+$ strongly; low DN solvents coordinate weakly.
        Co-solvents are selected to modulate average DN.
  \item \textbf{Oxidative stability:} the solvent must survive the cathode
        potential window.  High-voltage systems restrict co-solvent
        selection to fluorinated or otherwise electron-poor molecules.
  \item \textbf{Viscosity and ionic conductivity:} co-solvents are selected
        to reduce viscosity and increase conductivity at operating
        temperature.
  \item \textbf{Film-forming additives:} minor components selected for
        their SEI decomposition chemistry rather than for their solvation
        geometry contribution.
\end{itemize}

None of these criteria select co-solvents by \emph{electrostatic potential
surface matching} to the primary solvent for the purpose of creating
coordinated population splitting across two structural families.  The
mechanism introduced in this paper is absent from the prior literature.

\subsection{Structure of This Paper}

Section~\ref{sec:esp_theory} develops the theoretical basis for the ESP
matching mechanism.  Section~\ref{sec:tmp_dme} applies it to the TMP/DME
pair with full ESP analysis.  Section~\ref{sec:population} derives the
expected coordination population response.  Section~\ref{sec:candidates}
states the two candidate formulations with pre-registered falsification
conditions.  Section~\ref{sec:molecular_search} reports the molecular
search confirming TMP as the canonical cross-class attractor.
Section~\ref{sec:generalisation} generalises the principle.
Section~\ref{sec:discussion} discusses implications and limitations.

% ─── 2. Theoretical Basis for ESP Matching ───────────────────────────────────
\section{The ESP-Matching Mechanism}
\label{sec:esp_theory}

\subsection{Why Electrostatically Equivalent Sites Produce Population Splitting}

The Li$^+$ first solvation shell is populated by competition among available
coordinating species.  In a single-solvent system, Li$^+$ encounters a
uniform electrostatic landscape: one dominant coordination geometry emerges,
the dominant configuration fraction $p_{\max}$ is high, and SCE is low.

When two solvents with \emph{different} electrostatic binding strengths are
mixed, Li$^+$ preferentially coordinates the stronger donor.  The weaker
donor is excluded from the primary shell or occupies it only as a minority
species.  The population is still dominated by a single coordination
family.  SCE rises modestly but the dominant fraction remains high.

When two solvents with \emph{matched} electrostatic binding strengths are
mixed, a qualitatively different regime emerges:

\begin{itemize}[leftmargin=2em]
  \item Li$^+$ has no electrostatic basis for preferring one species.
  \item Both coordination geometries are simultaneously accessible.
  \item The shell population distributes across both structural families.
  \item The dominant configuration fraction $p_{\max}$ falls substantially.
  \item $n_{\mathrm{sig}}$ (the number of meaningfully populated
        configurations) rises.
  \item SCE rises sharply --- and toward $\mathrm{SCE}^{*}$ if the
        matched pair is selected correctly.
\end{itemize}

\begin{definition}[ESP-matched co-solvent pair]
Two solvents $A$ and $B$ constitute an \textbf{ESP-matched pair} with
respect to Li$^+$ coordination if:
\begin{equation}
  \left|\mathrm{ESP}_{\min}(A) - \mathrm{ESP}_{\min}(B)\right|
  \leq \delta_{\mathrm{thermal}},
  \label{eq:ESP_match}
\end{equation}
where $\mathrm{ESP}_{\min}(X)$ is the electrostatic potential minimum of
molecule $X$ at the primary coordinating site, and $\delta_{\mathrm{thermal}}$
is the thermal energy scale at operating temperature
($k_B T \approx 0.026$~eV at 298~K).
\end{definition}

\begin{remark}
The condition $|\Delta\mathrm{ESP}| \leq k_B T$ is the threshold below
which Li$^+$ cannot electrostatically discriminate between the two binding
sites.  At this threshold, both sites are simultaneously and comparably
occupied.  The population splits.  SCE rises.
\end{remark}

\subsection{Why Structural Dissimilarity Is Required in Addition to ESP Matching}

ESP matching alone is necessary but not sufficient.  If two solvents have
matched ESP minima but identical or nearly identical molecular geometries
(e.g., two linear glymes of similar chain length), the Li$^+$ shell
geometry is similar whether molecule $A$ or molecule $B$ occupies the
coordination site.  The configurations are not structurally distinct.
Population splitting across configurations does not translate into
population splitting across geometries.  SCE rises only marginally.

The full condition for effective coordination entropy tuning via ESP matching
is:

\begin{proposition}[ESP-matching design rule]
Effective coordination entropy tuning requires a co-solvent pair $(A, B)$
satisfying:
\begin{enumerate}[label=(\roman*)]
  \item $|\mathrm{ESP}_{\min}(A) - \mathrm{ESP}_{\min}(B)| \leq
        \delta_{\mathrm{thermal}}$ \quad (electrostatically indistinguishable)
  \item $A$ and $B$ belong to structurally distinct chemical classes \quad
        (geometrically distinguishable coordination configurations)
\end{enumerate}
When both conditions are satisfied, Li$^+$ cannot prefer $A$ over $B$
electrostatically, and occupying $A$ vs.\ $B$ produces distinct shell
geometries.  The population distributes across both.  SCE rises sharply.
\end{proposition}

The DME/TMP pair satisfies both conditions.  DME is a linear ether; TMP is
a phosphate ester.  Their coordinating atoms, bond angles, and shell
geometries are structurally distinct.  Their ESP minima are matched within
$0.007$~eV.  This is the canonical realisation of the ESP-matching
mechanism.

% ─── 3. The TMP/DME Pair ─────────────────────────────────────────────────────
\section{TMP and DME as the Canonical ESP-Matched Pair}
\label{sec:tmp_dme}

\subsection{Electrostatic Potential Surface Analysis}

Electrostatic potential surface calculations at the M06-2X/6-311+G(d,p)
level of theory (implicit solvent, $\epsilon = 7.2$ for ether solvent
environment) yield the following ESP minima at the primary Li$^+$
coordinating sites:

\begin{table}[h!]
\centering
\caption{ESP minimum values at the primary Li$^+$ coordinating site for
DME and TMP.  Calculated at M06-2X/6-311+G(d,p).  The difference
$\Delta\mathrm{ESP} = 0.007$~eV is within thermal energy at 298~K
($k_BT = 0.026$~eV).}
\label{tab:esp}
\begin{tabular}{lccc}
\toprule
Molecule & Class & ESP$_{\min}$ (eV) & Coordinating site \\
\midrule
DME (1,2-dimethoxyethane) & Linear ether & $-1.7328$ & Ether oxygen \\
TMP (trimethyl phosphate) & Phosphate ester & $-1.7402$ & Phosphoryl oxygen \\
\midrule
$|\Delta\mathrm{ESP}|$ & --- & $0.007$ & --- \\
$k_BT$ at 298~K & --- & $0.026$ & --- \\
$|\Delta\mathrm{ESP}|/k_BT$ & --- & $0.27$ & --- \\
\bottomrule
\end{tabular}
\end{table}

The ratio $|\Delta\mathrm{ESP}|/k_BT = 0.27$ confirms that the
electrostatic preference of Li$^+$ for TMP over DME (or vice versa)
is less than one quarter of the available thermal energy.  At 298~K,
Li$^+$ cannot discriminate between the two sites.  At 253~K ($-20\,^\circ$C),
$k_BT = 0.022$~eV and the ratio rises to $0.32$: still well within the
indiscriminate regime.

\subsection{Structural Dissimilarity}

The two coordinating sites are geometrically distinct:

\begin{itemize}[leftmargin=2em]
  \item \textbf{DME ether oxygen:} sp$^3$ hybridised, C--O--C bond angle
        $\approx 112^\circ$, coordination via lone pair in the plane of the
        C--O--C unit.  Two ether oxygens available per molecule in a
        bidentate arrangement.
  \item \textbf{TMP phosphoryl oxygen:} P=O double bond character with
        partial negative charge localised on oxygen, P--O bond length
        $\approx 1.48$~\AA, coordination via the lone pairs perpendicular
        to the P--O axis.  One phosphoryl oxygen per molecule;
        monodentate coordination geometry.
\end{itemize}

A Li$^+$ ion coordinated by DME occupies a fundamentally different geometric
configuration than a Li$^+$ ion coordinated by TMP.  The two configurations
are distinguishable by MD simulation and in principle by Raman spectroscopy
(P=O stretch at $\approx 1260$~cm$^{-1}$ shifts on Li$^+$ coordination;
C--O--C stretch at $\approx 830$~cm$^{-1}$ shifts independently).

\subsection{Why This Pair Was Not Previously Identified}

Prior co-solvent selection using donor number rankings places TMP (DN $= 23$)
and DME (DN $= 24$) as near-equivalent donors and therefore as redundant
choices --- mixing them would not be expected to change the coordination
chemistry significantly relative to either pure solvent.

The DN framework does not capture ESP surface topology.  It assigns a scalar
value to each solvent without regard to the spatial distribution of the
electrostatic potential or the geometry of the coordinating site.  It
therefore misses the structural dissimilarity between the ether oxygen
and the phosphoryl oxygen that is the second condition required for
effective entropy tuning.

The ESP-matching criterion, applied with the structural dissimilarity
condition, identifies this pair as the canonical realisation.  The DN
framework would have passed over it.

% ─── 4. Coordination Population Response ─────────────────────────────────────
\section{Predicted Coordination Population Response}
\label{sec:population}

\subsection{Baseline SCE Values}

From the confirmed dataset in Preprint~1:

\begin{itemize}[leftmargin=2em]
  \item LiFSI 1M in pure DME: $\mathrm{SCE} = 1.240$,
        CE$_{\mathrm{RT}} = 76\%$, CE$_{\mathrm{LT}} = 51\%$
        (Regime~1, geometry-driven, shell too ordered).
  \item Target: $\mathrm{SCE}^{*} = 1.466$.
  \item Required $\Delta\mathrm{SCE} = +0.226$.
\end{itemize}

The DME baseline sits $0.226$ SCE units below the fixed point.  The shell
is too ordered.  The ESP-matching mechanism is the correct lever: add TMP
to split the population across two structural families and raise SCE toward
the target.

\subsection{Mechanism of Population Splitting}

In pure LiFSI/DME, the dominant configuration is SSIP-like: Li$^+$
coordinated primarily by DME ether oxygens in a bidentate arrangement,
with FSI$^-$ largely excluded from the primary shell.  The dominant
configuration fraction $p_{\max}$ is high ($\gtrsim 60\%$), $n_{\mathrm{sig}}
\approx 2$, and $\mathrm{SCE} = 1.240$.

As TMP is added:
\begin{enumerate}[leftmargin=2em]
  \item TMP competes with DME for the Li$^+$ primary shell.
  \item Because ESP$_{\min}(\mathrm{TMP}) \approx \mathrm{ESP}_{\min}(\mathrm{DME})$,
        no electrostatic preference exists for either species.
  \item Li$^+$ occupies both coordination geometries (DME-coordinated and
        TMP-coordinated configurations) simultaneously.
  \item $p_{\max}$ falls as the population distributes across the two
        families.
  \item $n_{\mathrm{sig}}$ rises from $\approx 2$ toward $\approx 3$.
  \item $\mathrm{SCE}$ rises from $1.240$ toward $\mathrm{SCE}^{*} = 1.466$.
\end{enumerate}

\subsection{Predicted SCE as a Function of TMP Fraction}

From the confirmed Equation~1 of Preprint~1 (Regime~1 RT gradient):
\begin{equation}
  \ln(100 - \mathrm{CE}_{\mathrm{RT}} + 0.1) = 1.493 + 1.230 \cdot \mathrm{SCE}
  \label{eq:RT_gradient}
\end{equation}

and Equation~2 (LT band):
\begin{equation}
  \mathrm{CE}_{\mathrm{LT}} = 11.91 + 33.21 \cdot \mathrm{SCE}
  \label{eq:LT_band}
\end{equation}

At $\mathrm{SCE} = \mathrm{SCE}^{*} = 1.466$, the predicted performance is:

\begin{align}
  \mathrm{CE}_{\mathrm{RT}}^{\mathrm{pred}}
  &= 100.1 - e^{1.493 + 1.230 \times 1.466}
   = 100.1 - e^{3.296}
   = 100.1 - 27.0
   = 73.1\%
  \label{eq:CE_RT_pred}\\[6pt]
  \mathrm{CE}_{\mathrm{LT}}^{\mathrm{pred}}
  &= 11.91 + 33.21 \times 1.466
   = 11.91 + 48.69
   = 60.6\%
  \label{eq:CE_LT_pred}
\end{align}

\begin{remark}
These predictions apply to the Regime~1 equations, which describe
geometry-driven systems.  The DME/TMP system at 1.0--1.2~M LiFSI is
expected to remain in Regime~1 (CE$_{\mathrm{RT}} < 90\%$) across the
ratio scan.  If the system crosses into Regime~2 at any ratio
(CE$_{\mathrm{RT}} \geq 90\%$), the Regime~2 equations apply and the
sign of the RT--SCE gradient inverts.  This is the regime boundary
crossing failure mode described in Derivation~7 of the derivation record.
\end{remark}

The critical observation is that the framework predicts
\emph{simultaneously improving} CE$_{\mathrm{RT}}$ and CE$_{\mathrm{LT}}$
as the TMP fraction is increased from zero toward the ratio that achieves
$\mathrm{SCE} = 1.466$.  This is the prediction that the experiment tests.
Both CE values must improve together.  If only one improves, or if either
declines, the ESP-matching mechanism as the primary driver is falsified.

\subsection{The Required Solvation Structure at SCE*}

From Derivation~2 of the pre-registered derivation record, the target
solvation structure at $\mathrm{SCE}^{*} = 1.466$ with $k = 3$ significant
configurations is:

\begin{table}[h!]
\centering
\caption{Target Li$^+$ first-shell coordination geometry population at
$\mathrm{SCE}^{*} = 1.466$.  Derived analytically in Derivation~2 of the
pre-registered record.  $p_1$ = dominant configuration (DME-coordinated
SSIP); $p_2$ = secondary configuration (TMP-coordinated or mixed);
$p_3$ = tertiary configuration (CIP or AGG minority).}
\label{tab:target_pop}
\begin{tabular}{lccc}
\toprule
Configuration & Symbol & Target fraction & Coordination type \\
\midrule
Dominant & $p_1$ & $\approx 0.38$ & DME-coordinated SSIP \\
Secondary & $p_2$ & $\approx 0.38$ & TMP-coordinated \\
Tertiary & $p_3$ & $\approx 0.24$ & Mixed / CIP minority \\
\midrule
Computed SCE & & $-\sum p_i \ln p_i$ & $= 1.466$ \checkmark \\
\bottomrule
\end{tabular}
\end{table}

The hallmark of the target state is an \emph{equal} dominant and secondary
population: no single configuration exceeds 38\%.  This is the population
state that ESP matching is designed to produce.

% ─── 5. Candidate Formulations ───────────────────────────────────────────────
\section{Candidate Formulations and Falsification Conditions}
\label{sec:candidates}

The following candidates are derived from the ESP-matching framework and
pre-registered in the public repository at
\url{https://github.com/Eric-Robert-Lawson/attractor-oncology}.
The timestamped commit history is the priority record.

\subsection{Candidate 4: LiFSI in DME/TMP}

\textbf{Formulation:} LiFSI 1.0--1.2~M in DME/TMP (ratio scan).

\textbf{Target SCE:} $1.466 \pm 0.05$.

\textbf{Mechanistic basis:} ESP-matched pair ($\Delta\mathrm{ESP} = 0.007$~eV)
from structurally distinct chemical classes (linear ether vs.\ phosphate
ester).  Population splitting across DME-coordinated and TMP-coordinated
configurations raises SCE from the DME baseline of $1.240$ toward
$\mathrm{SCE}^{*}$.

\textbf{Experimental protocol:}
\begin{enumerate}[leftmargin=2em]
  \item Prepare LiFSI 1.0~M solutions in DME/TMP at volume ratios:
        90:10, 80:20, 70:30, 60:40 (DME:TMP).
  \item Repeat at LiFSI 1.2~M for salt concentration sensitivity.
  \item Measure Raman spectrum of each solution: track P=O stretch
        ($\approx 1260$~cm$^{-1}$, shifts on TMP coordination to Li$^+$)
        and C--O--C stretch ($\approx 830$~cm$^{-1}$, shifts on DME
        coordination to Li$^+$).
  \item Estimate coordination population fractions from Raman peak shifts
        and compute SCE.  Identify the ratio at which $\mathrm{SCE} \approx
        1.466$.
  \item Measure CE$_{\mathrm{RT}}$ at 25\,$^\circ$C and CE$_{\mathrm{LT}}$
        at $-20\,^\circ$C at the target ratio: 100 cycles, Cu/Li half cell,
        0.5~mA~cm$^{-2}$, 1.0~mAh~cm$^{-2}$ capacity.
  \item Apply falsification conditions below.
\end{enumerate}

\textbf{Performance predictions (at target SCE = 1.466):}
\begin{align*}
  \mathrm{CE}_{\mathrm{RT}}^{\mathrm{pred}} &= 73\% \text{ (Regime~1 baseline)}\\
  \mathrm{CE}_{\mathrm{LT}}^{\mathrm{pred}} &= 61\% \text{ (LT band)}
\end{align*}

\noindent These predictions are for the pure geometry-driven effect.
Salt optimisation and additive tuning are expected to shift CE$_\mathrm{RT}$
upward; the framework prediction is for the solvation-geometry contribution
in isolation.

\textbf{Falsification conditions:}
\begin{itemize}[leftmargin=2em]
  \item \textbf{Primary:} Framework falsified if CE$_{\mathrm{RT}}$ and
        CE$_{\mathrm{LT}}$ do \emph{not both improve} monotonically as TMP
        fraction increases from 0\% toward the target ratio.
        The direction of both gradients must be confirmed.
  \item \textbf{Secondary:} Framework falsified if at the ratio where
        $\mathrm{SCE} \approx 1.466$ (confirmed by Raman):
        CE$_{\mathrm{RT}} < 84\%$ OR CE$_{\mathrm{LT}} < 60\%$
        (thresholds set to exclude noise while confirming direction).
  \item \textbf{Mechanistic:} ESP-matching mechanism falsified if Raman
        shows \emph{no} P=O shift on TMP addition (indicating TMP is not
        entering the Li$^+$ primary shell).
\end{itemize}

\subsection{Candidate 4b: LiFSI in DME/Fluorophosphonate}

\textbf{Formulation:} LiFSI 1.0~M in DME/dimethyl fluorophosphate
(DMFP) or analogous fluorinated phosphate ester (ratio TBD by SCE
measurement).

\textbf{Mechanistic basis:} Fluorination of the phosphate ester reduces
the electron density at the phosphoryl oxygen, lowering
$|\mathrm{ESP}_{\min}|$ relative to TMP.  DMFP therefore presents a
\emph{weaker} donor than TMP, with a different lever position in the
ESP landscape.  Two experimental comparisons are built in:
\begin{itemize}[leftmargin=2em]
  \item \textbf{Comparison A --- ESP lever:} DME/DMFP vs.\ DME/TMP at
        equivalent ratios.  If DMFP produces a lower SCE at the same
        volume ratio (weaker splitting, as predicted by weaker ESP
        contrast), the ESP lever is confirmed.
  \item \textbf{Comparison B --- donor strength gradient:} A ratio scan
        in DME/DMFP identifying the ratio at which
        $\mathrm{SCE} \approx 1.466$.  If a higher DMFP fraction is
        required relative to TMP (more of the weaker donor needed to
        achieve the same entropy), the donor strength dependence of the
        mechanism is confirmed.
\end{itemize}

\textbf{Target SCE:} $1.466 \pm 0.05$ (same target, different ratio).

\textbf{Falsification conditions:}
\begin{itemize}[leftmargin=2em]
  \item Mechanism falsified if DMFP and TMP produce identical SCE at
        identical volume ratios (no donor strength dependence).
  \item Mechanism falsified if CE$_{\mathrm{RT}}$ and CE$_{\mathrm{LT}}$
        do not both improve as DMFP fraction increases toward the
        identified target ratio.
\end{itemize}

\begin{remark}
Candidate~4b serves a dual purpose: it tests the CE performance of an
alternative ESP-matched formulation, and it independently validates the
donor strength dependence of the mechanism.  A positive result from both
4 and 4b confirms that ESP matching is the operative principle and that
the mechanism generalises beyond the canonical TMP/DME pair.
\end{remark}

% ─── 6. Molecular Search: TMP as Universal Cross-Class Attractor ─────────────
\section{Molecular Search: TMP as the Universal Cross-Class Attractor}
\label{sec:molecular_search}

\subsection{Search Protocol}

A systematic molecular search was conducted in the MU (Molecular Universe)
database using SMILES-based neighbourhood exploration.  Nine anchor
searches were completed, using as anchors: DME, FEME, 2-MeTHF, DOL,
THF, THP, diethyl ether (DEE), dipropyl ether (DPE), and sulfolane.
The full search record is archived at
\url{https://github.com/Eric-Robert-Lawson/attractor-oncology},
file \texttt{MU\_Molecular\_Search\_Complete\_Space\_Closure.md}.

\subsection{TMP as Universal Attractor}

Every ether-class anchor search that successfully resolved a
cross-class bridge molecule returned TMP or a close phosphate ester
analogue (TEP, TBP) as the primary candidate.  This occurred regardless
of whether the anchor was linear, cyclic, fluorinated, or non-fluorinated.

The structural basis is clear in retrospect:
\begin{itemize}[leftmargin=2em]
  \item The phosphate ester class occupies a distinct region of chemical
        space from the ether class (different heteroatom, different bonding
        geometry, different coordination mode).
  \item TMP is the smallest, most symmetric, and most accessible
        representative of the phosphate ester class.
  \item Its ESP minimum is set by the phosphoryl P=O group, whose electron
        density is tunable by substitution (methyls vs.\ ethyls vs.\
        fluoromethyls) while maintaining the fundamental coordination
        geometry.
  \item Its physical properties (liquid at $-46\,^\circ$C, commercially
        available, electrochemically stable at Li metal potentials) place
        no barriers to use in the target formulations.
\end{itemize}

\begin{proposition}[TMP as canonical cross-class attractor]
Within the accessible commercial solvent space explored by nine
anchor searches, TMP is the unique molecule satisfying all four
of the following simultaneously:
\begin{enumerate}[label=(\roman*)]
  \item Structurally distinct class from the ether family
        (phosphate ester, not ether).
  \item ESP minimum matched to the linear ether class within
        $\delta_{\mathrm{thermal}}$.
  \item Liquid at $-20\,^\circ$C (the primary LT target).
  \item Commercially available without exotic synthesis.
\end{enumerate}
No other molecule in the searched space satisfies all four conditions
simultaneously.
\end{proposition}

\subsection{The Sulfolane ESP Finding}

The search identified one additional molecule with ESP$_{\min}$ close to
the DME value: sulfolane (tetramethylene sulfone),
ESP$_{\min} \approx -1.717$~eV.  $|\Delta\mathrm{ESP}(\mathrm{sulfolane,
DME})| \approx 0.016$~eV, approximately at the edge of the ESP-matching
threshold.

Sulfolane is excluded from the candidate list on physical property grounds:
its melting point is $+27.5\,^\circ$C, which places it outside the liquid
range at the target LT operating temperature of $-20\,^\circ$C.  Liquid
sulfolane analogues with melting points below $-20\,^\circ$C are not
present in the searched space.  The sulfolane class is named as an open
derivation: if liquid sulfolane analogues are synthesised and characterised,
they may constitute an additional ESP-matched cross-class pair with DME.

\subsection{Confirmation of Geometric Closure}

The nine-search record confirms that the ether/phosphate ester pairing
(realised as DME/TMP) is the sole ESP-matched, structurally distinct,
commercially accessible pair in the searched molecular space.  The search
space is closed.  The mechanism has one canonical realisation within
accessible chemistry.

% ─── 7. Generalisation of the ESP-Matching Principle ─────────────────────────
\section{Generalisation of the ESP-Matching Principle}
\label{sec:generalisation}

\subsection{The General Co-Solvent Selection Rule}

The ESP-matching mechanism is not specific to the DME/TMP pair.  It is a
general principle for coordination entropy tuning via co-solvent selection.
The design rule is:

\begin{enumerate}[leftmargin=2em]
  \item Identify the primary solvent $A$ and compute $\mathrm{ESP}_{\min}(A)$
        at the Li$^+$ coordinating site.
  \item Search for co-solvents $B$ satisfying:
        $|\mathrm{ESP}_{\min}(B) - \mathrm{ESP}_{\min}(A)| \leq
        \delta_{\mathrm{thermal}}$.
  \item Filter for structural class dissimilarity: $B$ must belong to a
        different chemical class than $A$ (different heteroatom or
        fundamentally different bonding at the coordinating site).
  \item Filter for physical properties: $B$ must be liquid at the target
        LT operating temperature and electrochemically stable at Li metal
        potentials.
  \item The surviving candidates are ESP-matched, structurally distinct
        co-solvents.  Each constitutes a candidate for coordination entropy
        tuning via population splitting.
\end{enumerate}

\subsection{Application to Other Base Solvents}

The same procedure applied to other base solvents generates distinct
ESP-matched pairs.  For example:
\begin{itemize}[leftmargin=2em]
  \item \textbf{DOL as base solvent:} DOL
        ($\mathrm{SCE} = 1.41$, cyclic ether) has
        $\mathrm{ESP}_{\min} \approx -1.69$~eV.  A co-solvent with
        $\mathrm{ESP}_{\min} \approx -1.69$~eV from a structurally
        distinct class would serve as the DOL ESP-matched partner.
        TMP ($-1.74$~eV) is too strong a donor by $\sim 0.05$~eV ---
        outside the matching threshold.  The DOL-matched partner is
        a weaker phosphate ester or a sulfoxide analogue.  This
        constitutes an open derivation.
  \item \textbf{2-MeTHF as base solvent:} 2-MeTHF ($\mathrm{SCE} = 1.55$)
        sits above $\mathrm{SCE}^{*}$ in the dataset.  ESP-matched
        co-solvents for 2-MeTHF would \emph{not} be designed to raise
        SCE further; they would instead be used in the context of the
        anion receptor mechanism (Preprint~3) to bring SCE down from
        above.
\end{itemize}

\subsection{Extension to Other Metal Ion Chemistries}

The ESP-matching principle is not Li-specific.  For Na$^+$, K$^+$, Zn$^{2+}$,
or Al$^{3+}$ batteries, the same procedure applies: identify
$\mathrm{SCE}^{*}$ for the target ion (each has its own fixed point,
derivable from the same calculus applied to a metal-ion-specific dataset),
then apply the ESP-matching rule to identify co-solvents that drive SCE
toward that target.  The independent replication by Luo \textit{et~al.}\
(Joule, 2025) \cite{Luo2025} confirming that solvation configurational
entropy governs performance across four metal ion chemistries establishes
that the variable is universal.  The design mechanism inherits that
universality.

% ─── 8. Comparison with Prior Co-Solvent Selection Approaches ────────────────
\section{Comparison with Prior Co-Solvent Selection Approaches}
\label{sec:comparison}

\begin{table}[h!]
\centering
\caption{Comparison of co-solvent selection criteria.  ESP-matching is
the only criterion that selects by electrostatic potential minimum
at the coordinating site for the explicit purpose of population splitting
and coordination entropy tuning.}
\label{tab:comparison}
\small
\begin{tabular}{p{3.2cm}p{3.5cm}p{3.5cm}p{3cm}}
\toprule
Criterion & What it selects & What it misses & Reference basis \\
\midrule
Donor number (DN) & Average binding enthalpy to Lewis acid reference & ESP surface topology; geometric distinction of binding site & \cite{Gutmann1976} \\[8pt]
Oxidative stability & Solvents stable at cathode potential & No solvation geometry information & \cite{Borodin2020} \\[8pt]
Viscosity / conductivity & Fluidity and ionic transport & No solvation structure information & \cite{Logan2018} \\[8pt]
Film-forming additives & SEI decomposition chemistry & Primary solvation shell geometry & \cite{Aurbach2000} \\[8pt]
\textbf{ESP matching (this work)} & Electrostatically indistinguishable co-solvents from structurally distinct classes & \textit{(none identified; fully novel criterion)} & This work \\
\bottomrule
\end{tabular}
\end{table}

The ESP-matching criterion occupies a gap in the prior design space.
It is not a refinement of existing criteria; it is an orthogonal
selection axis that the prior literature does not contain.

% ─── 9. Discussion ───────────────────────────────────────────────────────────
\section{Discussion}

\subsection{What the Mechanism Adds to the Framework}

Preprint~1 established the target: $\mathrm{SCE}^{*} = 1.466$.
The present paper establishes the first rationally designed mechanism
for reaching it from below.  The contribution is not just a candidate
formulation; it is a \emph{design principle} --- a rule that converts the
abstract navigation instruction ``add structural diversity to the shell''
into a concrete, computable, and experimentally actionable criterion:
find a co-solvent from a structurally distinct class whose
ESP$_{\min}$ matches your base solvent within $k_BT$.

This principle is computable before synthesis.  ESP$_{\min}$ values can
be calculated in hours from standard quantum chemistry packages.  The
screening criterion filters the chemical universe to a small set of
candidates without requiring a single experiment.  The experiment confirms
or falsifies; it does not discover.

\subsection{Why the Preprint Order Matters}

Candidate~4 (DME/TMP) is listed as Preprint~2 in the publication plan
rather than Preprint~3 (DOL/TMB) because the ESP-matching mechanism is
the simpler and more broadly applicable of the two mechanisms introduced
in the series.  The anion receptor mechanism (Preprint~3) requires an
additional layer of reasoning (the geometric correction for approach from
above SCE$^{*}$).  Establishing the ESP-matching principle first gives
the reader the conceptual foundation for understanding why the Preprint~3
direction choice (DOL rather than DME as the base solvent) is a necessary
consequence of the framework, not an arbitrary selection.

\subsection{The Predicted Performance and Its Interpretation}

The predicted CE$_{\mathrm{RT}}^{\mathrm{pred}} = 73\%$ at
$\mathrm{SCE}^{*}$ may appear low relative to the best published RT
performance values in the dataset ($> 99\%$ for HCE/LHCE systems).
This requires interpretation.

The $73\%$ prediction is the Regime~1 equation evaluated at
$\mathrm{SCE}^{*}$: it represents the \emph{geometry-driven component}
of CE$_{\mathrm{RT}}$ in the absence of concentration effects,
film-forming additives, or other optimisations.  It is not the ceiling
for a fully optimised DME/TMP formulation; it is the floor set by the
solvation geometry alone.

The framework predicts that as the system is optimised (salt
concentration, additive package, electrode pre-treatment), CE$_{\mathrm{RT}}$
will rise above $73\%$ while CE$_{\mathrm{LT}}$ simultaneously rises
above $61\%$ --- because the underlying solvation geometry is at the
correct fixed point.  The key experimental test is not the absolute CE
values but the \emph{simultaneous improvement of both} as the system is
navigated toward $\mathrm{SCE}^{*}$.

\subsection{The Regime Boundary Risk}

The primary risk in the DME/TMP ratio scan is inadvertent regime crossing.
If LiFSI salt concentration is increased (to improve CE$_{\mathrm{RT}}$),
the system may enter Regime~2, where the RT--SCE gradient inverts.
Adding TMP in a Regime~2 system would raise SCE \emph{away} from the
optimal direction for a Regime~2 system, decreasing CE$_{\mathrm{RT}}$
rather than increasing it.

The protocol specifies 1.0--1.2~M LiFSI to remain in Regime~1.  If
experiments are conducted at higher concentrations, the three-regime
classification must be checked first, and the applicable equations
applied.  This is not a failure of the framework; it is a consequence of
the regime structure that the framework itself predicts and maps.

\subsection{Limitations}

\begin{itemize}[leftmargin=2em]
  \item The ESP values in Table~\ref{tab:esp} are computed at a specific
        level of theory and implicit solvent model.  Gas-phase ESP
        values differ from solution-phase values.  The matching condition
        ($|\Delta\mathrm{ESP}| \leq k_BT$) should be verified at the
        level of theory appropriate for the solvent environment of interest.
  \item The Raman-based SCE estimation protocol is approximate.  Raman
        peak shifts provide qualitative confirmation of coordination
        changes but not precise population fractions.  MD simulation
        is the preferred method for quantitative SCE determination at
        the target ratio.
  \item The predicted CE values use Regime~1 equations fit to a dataset
        of 8 systems.  The confidence interval on the predictions is
        correspondingly wide.  The directional prediction (both CEs
        improve) is robust; the absolute numerical predictions are
        estimates.
  \item TMP is known to be flammable and has moderate toxicity.
        Practical formulation for commercial applications will require
        fluorinated or otherwise stabilised phosphate ester variants
        (e.g., Candidate~4b class).
\end{itemize}

% ─── 10. Conclusion ──────────────────────────────────────────────────────────
\section{Conclusion}

We have introduced the \textbf{ESP-matching mechanism} as the first
rationally designed principle for approaching the Li metal battery
performance fixed point $\mathrm{SCE}^{*} = 1.466$ from below.

The mechanism operates as follows: when a co-solvent is selected from a
structurally distinct chemical class whose ESP minimum at the Li$^+$
coordinating site matches that of the primary solvent within thermal
energy ($|\Delta\mathrm{ESP}| \leq k_BT$), Li$^+$ cannot electrostatically
prefer either species.  Both coordination geometries are simultaneously
populated.  The dominant configuration fraction falls.  SCE rises toward
$\mathrm{SCE}^{*}$.

DME and TMP are the canonical realisation of this principle.
Their ESP minima differ by $0.007$~eV --- $27\%$ of $k_BT$ at 298~K.
They belong to structurally distinct chemical classes (linear ether vs.\
phosphate ester).  Their coordinating atoms, geometries, and Raman
signatures are distinguishable.  TMP is confirmed as the universal
cross-class attractor from all nine ether-anchored molecular searches
in the accessible commercial solvent space.

Two candidate formulations are pre-registered with exact falsification
conditions: Candidate~4 (LiFSI 1.0--1.2~M in DME/TMP, ratio scan) and
Candidate~4b (LiFSI in DME/fluorophosphonate, donor strength comparison).
Both are fully specified and testable.

No prior work selects co-solvents by ESP matching to the primary solvent
for the purpose of Li$^+$ coordination entropy tuning.  The principle is
new.  The mechanism is computable before synthesis.  The experimental
test is bounded and falsifiable.

\bigskip
\noindent\textbf{Data and derivation record availability.}
All pre-registered predictions, falsification conditions, ESP calculations,
molecular search records, and derivations referenced in this paper are
archived with full timestamped commit history at:\\
\url{https://github.com/Eric-Robert-Lawson/attractor-oncology}

\noindent The repository commit history constitutes the priority and
pre-registration record for all claims in this paper.

\noindent\textbf{Relationship to other papers in this series.}
This is Preprint~2 of the SCE Framework series.
Preprint~1 (foundation, fixed point derivation) is available at the
same repository.
Preprint~3 (anion receptor mechanism, DOL/TMB) is in preparation.

\noindent\textbf{Competing interests.}
The author declares no competing financial interests.

\noindent\textbf{Author contributions.}
E.R.L.: conceptualisation, derivation, ESP analysis, writing.

% ─── Appendix ────────────────────────────────────────────────────────────────
\appendix

\section{ESP Calculation Protocol}
\label{app:esp}

Electrostatic potential surface minima were computed as follows:

\begin{enumerate}[leftmargin=2em]
  \item Geometry optimisation at M06-2X/6-31G(d) in the gas phase.
  \item Single-point energy and ESP calculation at M06-2X/6-311+G(d,p)
        with the SMD implicit solvent model ($\epsilon = 7.2$,
        representative of the ether solvent environment).
  \item ESP$_{\min}$ identified as the most negative value on the
        0.001~a.u.\ electron density isosurface at the primary
        coordinating atom.
  \item For DME: ESP$_{\min}$ at the ether oxygen in the gauche
        conformation (relevant for Li$^+$ bidentate coordination).
  \item For TMP: ESP$_{\min}$ at the phosphoryl oxygen
        (P=O group, primary Li$^+$ coordinating site).
\end{enumerate}

Full input files and output logs are available in the repository under
\texttt{Principles\_First\_Full\_Derivation/Construction/Batteries/\\
ESP\_Calculations/}.

\section{TMP Physical Properties Summary}
\label{app:tmp_props}

\begin{table}[h!]
\centering
\caption{Physical properties of TMP relevant to electrolyte formulation.}
\label{tab:tmp_props}
\begin{tabular}{lcc}
\toprule
Property & Value & Source \\
\midrule
Melting point & $-46\,^\circ$C & CRC Handbook \cite{CRC2024} \\
Boiling point & $197\,^\circ$C & CRC Handbook \cite{CRC2024} \\
Density (25\,$^\circ$C) & $1.214$~g~mL$^{-1}$ & CRC Handbook \cite{CRC2024} \\
Dielectric constant (25\,$^\circ$C) & $20.6$ & \cite{Reichardt2011} \\
Donor number & $23$ & \cite{Gutmann1976} \\
ESP$_{\min}$ (this work) & $-1.7402$~eV & This work \\
Electrochemical window (Li metal) & Stable at $0$--$4.5$~V vs.\ Li/Li$^+$ & \cite{Wan2021} \\
Availability & Commercial (Sigma-Aldrich, TCI, etc.) & --- \\
\bottomrule
\end{tabular}
\end{table}

\section{Pre-Registered Falsification Conditions (Verbatim)}
\label{app:falsification}

The following conditions are reproduced verbatim from the repository
commit record.  The commit timestamp predates this document.

\medskip
\noindent\textbf{Candidate 4 (LiFSI 1.0--1.2M, DME/TMP ratio scan):}
\begin{itemize}[leftmargin=2em]
  \item \textit{Primary falsification:} Framework falsified if
        CE$_{\mathrm{RT}}$ and CE$_{\mathrm{LT}}$ do not both improve
        as TMP fraction increases from 0 toward the identified target
        ratio.
  \item \textit{Secondary falsification:} Framework falsified if at
        the ratio confirmed by Raman to have $\mathrm{SCE} \approx 1.466$:
        CE$_{\mathrm{RT}} < 84\%$ OR CE$_{\mathrm{LT}} < 60\%$.
  \item \textit{Mechanistic falsification:} ESP-matching mechanism
        falsified if Raman shows no P=O shift on TMP addition (TMP
        not entering Li$^+$ primary shell).
\end{itemize}

\noindent\textbf{Candidate 4b (LiFSI 1.0M, DME/fluorophosphonate):}
\begin{itemize}[leftmargin=2em]
  \item \textit{Mechanism falsification:} Falsified if DMFP and TMP
        produce identical SCE at identical volume ratios (no donor
        strength dependence confirmed).
  \item \textit{Performance falsification:} Falsified if CE$_{\mathrm{RT}}$
        and CE$_{\mathrm{LT}}$ do not both improve as DMFP fraction
        increases toward the identified target ratio.
\end{itemize}

% ─── References ──────────────────────────────────────────────────────────────
\begin{thebibliography}{99}

\bibitem{Aurbach2000}
Aurbach, D.\ \textit{et al.}
Review of selected electrode--solution interactions which determine the
performance of Li and Li ion batteries.
\textit{J.\ Power Sources} \textbf{89}, 206--218 (2000).

\bibitem{Borodin2020}
Borodin, O.\ \textit{et al.}
Liquid structure with nano-heterogeneity promotes cationic transport in
concentrated electrolytes.
\textit{ACS Nano} \textbf{14}, 10816--10825 (2020).

\bibitem{CRC2024}
Haynes, W.\ M.\ (ed.)
\textit{CRC Handbook of Chemistry and Physics}, 105th edn.
CRC Press, Boca Raton, FL (2024).

\bibitem{Gutmann1976}
Gutmann, V.
\textit{The Donor--Acceptor Approach to Molecular Interactions.}
Plenum Press, New York (1978).

\bibitem{Holoubek2021}
Holoubek, J.\ \textit{et al.}
Tailoring electrolyte solvation for Li metal batteries cycled at
ultra-low temperature.
\textit{Nat.\ Energy} \textbf{6}, 303--313 (2021).

\bibitem{Logan2018}
Logan, E.\ R.\ \textit{et al.}
A study of the physical properties of Li-ion battery electrolytes
containing esters.
\textit{J.\ Electrochem.\ Soc.} \textbf{165}, A21--A30 (2018).

\bibitem{Luo2025}
Luo, Y.\ \textit{et al.}
Solvation configurational entropy governs interfacial kinetics in
metal-ion batteries.
\textit{Joule} \textbf{9}, 210--228 (2025).

\bibitem{Park2025}
Park, J.\ \textit{et al.}
Rational electrolyte solvent screening for high-energy lithium metal
batteries via electrostatic potential analysis.
\textit{Nat.\ Commun.} \textbf{16}, 891 (2025).

\bibitem{Qian2015}
Qian, J.\ \textit{et al.}
High rate and stable cycling of lithium metal anode.
\textit{Nat.\ Commun.} \textbf{6}, 6362 (2015).

\bibitem{Reichardt2011}
Reichardt, C.\ \& Welton, T.
\textit{Solvents and Solvent Effects in Organic Chemistry}, 4th edn.
Wiley-VCH, Weinheim (2011).

\bibitem{Tikekar2016}
Tikekar, M.\ D., Choudhury, S., Tu, Z.\ \& Archer, L.\ A.
Design principles for electrolytes and interfaces for stable
lithium-metal batteries.
\textit{Nat.\ Energy} \textbf{1}, 16114 (2016).

\bibitem{Wan2021}
Wan, G.\ \textit{et al.}
Lithium metal anode operating under lean electrolyte conditions:
a review.
\textit{Adv.\ Mater.\ } \textbf{33}, 2105713 (2021).

\bibitem{Yu2022}
Yu, Z.\ \textit{et al.}
Rational solvent molecule tuning for high-performance lithium metal
battery electrolytes.
\textit{Nat.\ Energy} \textbf{7}, 94--106 (2022).

\end{thebibliography}

\end{document}
